\begin{abstract}
The persistent challenge of affordable and responsive access to space for nanosatellites motivates examination of unconventional launch architectures. This paper investigates a modular stratospheric platform based on hybrid aerostats operating at equatorial latitudes to support a horizontal pneumatic accelerator for nanosatellite launch. We analyze the concept's potential to address limitations of conventional rockets (cost, responsiveness, infrastructure), circumvent space elevator material constraints, and complement electromagnetic launchers and mass drivers by leveraging atmospheric properties at altitude. Following a critical review of high-altitude launch precedents, we identify a specific research gap: the integration of stratospheric station-keeping platforms with kilometer-scale horizontal evacuated tubes for hypersonic pneumatic acceleration. The central research question is formulated: Can a modular stratospheric aerostat platform support a practical horizontal pneumatic launch system for nanosatellites? We conclude that feasibility remains unproven and outline a rigorous, phased research approach prioritizing physics-based modeling and subscale validation before proceeding to integrated system tests. All numerical parameters are treated as hypotheses requiring verification, with careful distinction between established physics, engineering estimates, proposed values, and speculation.
\end{abstract}
